\section{Conclusion}%
\label{sec:conclusion}
In this work, we explore generative self-supervised learning in graphs and identify 
the common issues that are faced by graph autoencoders. 
We present \model---a simple masked graph autoencoder---to address them from the  perspective of reconstruction objective, learning, loss function, and model architecture.
%the challenges of current graph self-supervised learning. 
In \model, we design the masked feature reconstruction strategy with a scaled cosine error as the reconstruction criterion.  
%from the perspective of objective, learning, loss function, and model architecture.
We conduct extensive experiments on a wide range of node and graph classification benchmarks, and the results demonstrate the  effectiveness and generalizability of \model. 
Our work shows that generative SSL can offer great potential to graph representation learning and pre-training, requiring  
%and deserve 
more in-depth explorations for future work.
% \newpage





\hide{


\section{Conclusion}%
\label{sec:conclusion}
In this work, we explore generative learning in graphs and present the challenges of current graph self-supervised learning. Then we propose \model a masked graph autoencoder, to address the problems from the perspective of objective, learning, loss function, and model architecture.
We conduct extensive experiments on a wide range of node and graph classification benchmarks, and the results demonstrate the generalizability and effectiveness of \model. 
Our work shows that generative SSL can offer great potential to graph representation learning and deserve more in-depth exploration in the future.
% \newpage



}